% Problemas del capítulo 15. Generado por tools/problemas/gen.py; no editar a mano.
\section{Problemas}\label{sec:problemas-15}

Graduados: \({\star}\) directo, \({\star}{\star}\) para pensar, \({\star}{\star}{\star}\) desafío. Las respuestas están en el apéndice~\ref{ap:respuestas}; conviene intentarlos antes de mirarlas.

\begin{problema}[\({\star}\)]\label{pb:15-1}
En una esfera de radio \(6371\)~km, ¿cuánto mide el arco del ecuador entre las longitudes \(0^\circ\) y \(90^\circ\)?
\end{problema}

\begin{problema}[\({\star}{\star}\)]\label{pb:15-2}
Un triángulo tiene un vértice en el polo y dos en el ecuador, separados \(60^\circ\) de longitud. ¿Cuánto valen sus tres ángulos? Calculá su exceso esférico y su área. ¿Qué fracción de la esfera es?
\end{problema}

\begin{problema}[\({\star}{\star}\)]\label{pb:15-3}
Tomando la Tierra como esfera de \(6371\)~km, calculá la distancia por el camino más corto entre Buenos Aires (\(-34{,}6^\circ\); \(-58{,}4^\circ\)) y Madrid (\(40{,}4^\circ\); \(-3{,}7^\circ\)), y el acimut con el que hay que salir.
\end{problema}

\begin{problema}[\({\star}{\star}\)]\label{pb:15-4}
Un triángulo rectángulo dibujado en una cancha tiene catetos de \(100\) y \(70\)~m. ¿Cuál es su exceso esférico, en segundos de arco?
\end{problema}

\begin{problema}[\({\star}{\star}{\star}\)]\label{pb:15-5}
En la proyección de Mercator, el factor de escala en la latitud \(\varphi\) es \(1/\cos\varphi\). ¿Cuánto se agrandan las distancias y las áreas a \(72^\circ\) de latitud, en el centro de Groenlandia, respecto del ecuador?
\end{problema}
